\uuid{LTv9}
\exo7id{4301}
\titre{exo7 4301}
\auteur{quercia}
\organisation{exo7}
\datecreate{2010-03-12}
\isIndication{false}
\isCorrection{true}
\chapitre{Calcul d'intégrales}
\sousChapitre{Intégrale impropre}
\cpge{mpsi-mp2i.analyse-asymptotique.a, pcsi.analyse-asymptotique.a, ptsi.analyse-asymptotique.a, mpsi-mp2i.integration.c, pcsi.integration.b, ptsi.integration.b, bcpst1.limites-continuite.b, bcpst1.integration-segment.d, mp-mpi.integration-intervalle.b, bcpst2.integrales-generalisees}
\module{Analyse}
\niveau{L1}
\difficulte{}

\contenu{
\texte{
On pose $f(x) =  \int_{t=x}^{+\infty} \frac{e^{-t}}t \,d t$.
}
\begin{enumerate}
    \item \question{Chercher $\lim_{x\to{+\infty}} f(x)$.}
    \item \question{A l'aide d'une intégration par parties, donner un équivalent de $f(x)$
      pour $x\to{+\infty}$.}
    \item \question{Donner un équivalent de $f(x)$ pour $x\to0^+$.}
\reponse{
$e^{-x}/x$.
$-\ln x$.
}
\end{enumerate}
}
